\section{Conclusion and Outlook}
\label{sec:conclusion}
%------------------------------------------------------------------------------

Since Feynman proposed the parton model more than fifty years ago, our understanding of the partonic structure of the proton has been greatly advanced. On one hand, a number of high-energy experiments carried out at facilities worldwide including SLAC, DESY, CERN, Fermi Lab, JLab, BNL, etc. allowed us to probe various aspects of hadronic structures at different energies and polarizations. On the other hand, many parton observables have been proposed in parallel
that provide a multi-dimensional description of the proton structure, including the collinear
PDFs, TMDPDFs, GPDs, parton DAs, LFWFs and so on.

Although QCD factorization theorems with RG improvement allow us to extract these parton observables through their connection to experimental observables, it is highly desirable to predict them from {\it ab initio}
calculations such as lattice QCD. Developments along this line have been rather slow due to
difficulties in simulating real-time dynamics. The situation, however, has changed since the proposal of LaMET a few years ago, which provides a systematically improvable method to calculate parton physics from first principles.

In this paper, we give an overview of LaMET formalism and its applications to observables which
can be accessed in lattice QCD and other non-perturbative methods.  By investigating the frame dependence of the structure of bound state hadrons, we explain how the IMF physics or parton physics naturally arises as an EFT description of the proton structure. Such an EFT description is most naturally formulated in SCET and LFQ,
but practical non-perturbative calculations of the proton matrix elements
have been difficult. LaMET in effect provides what is needed to realize LFQ. This is achieved by forming appropriate quasi parton observables in a large momentum state and match them to the true parton observables on the LF through factorization. In the case of PDFs, the former corresponds to finite-momentum distributions whose running is controlled by the momentum RGE, whereas the latter corresponds to IMF PDFs whose running is controlled by the usual RGE. It should be pointed out that
LaMET is a very general framework which can be applied to
large-momentum physical quantities calculated with any non-perturbative methods,
either Euclidean (with imaginary time) or Minkowskian (with real time).
Moreover, given a large momentum state, the same parton physics can be determined
from different quasi observables that form a universality class.

We then present how to calculate the parton observables in practice, with a particular focus on the collinear PDFs, GPDs, DAs, TMDPDFs and LFWFs. We also discuss the proton spin structure and show how the partonic contributions to proton spin can be obtained following the same approach. We finally summarize the lattice studies carried out so far with LaMET which, on one hand, demonstrate that LaMET is a promising approach to compute partonic structures of the proton, and on the other hand, clearly indicate that a lot of improvements are still required to reach such an accuracy that the lattice results can have considerable impact on phenomenology.

We complete this review with a few comments on improvements of lattice calculations for the future.
%
We recommend~\cite{Alexandrou:2019lfo} for more systematic discussion on some of the issues, for example, the continuum, infinite volume, and physical pion mass limits.

\begin{itemize}
	
	\item Large hadron momentum. Since the future of LaMET lies in larger momenta which naturally require smaller lattice spacings, it will be critical to address the challenges from using large momenta and small spacings for exa-scale computations, such as the excited state contamination or topological charge freezing problem.
	
	\item Renormalization. As discussed in \sec{renorm-npr}, the mass renormalization of Wilson line operators is favored for it is gauge invariant and does not introduce extra higher-twist effects or large statistical errors at long distance. However, its matching to the $\MS$ scheme, especially the renormalon ambiguities, still needs to be resolved for a full systematic application. Moreover, alternative schemes that include the above features are also highly desirable.
	
	\item Higher-order perturbative matching. In current LaMET calculations, one-loop perturbative matching
	has brought considerable corrections. Higher-order matching kernels will be necessary to control the systematics from this procedure.
	
	\item Power corrections. They are important if the hadron momentum is not very large or when $x$ is close to 0 or 1.
	Little progress has been made toward a model-independent determination of the power corrections so far. One contingent strategy is to extrapolate to $P^z\to\infty$ limit after implementing matching and target-mass corrections, but the ultimate solution relies on the lattice calculation of higher-twist distributions that has been discussed in \sec{others-ht}.
	
\end{itemize}

The above discussion of systematics is generic and applies to all quasi-observables. The rich theoretical developments in the past years have paved the way for calculating a wide range of parton observables using LaMET. With the rapid increase in computing resources and progress in developing new techniques and algorithms, we expect to see the above systematics to be kept under control step by step in the future. That would be important in establishing LaMET as a systematic approach to computing parton physics, and making lattice calculations play a crucial role in the EIC era.


%------------------------------------------------------------------------------
